De overwinning klonk.
Zijn visserssectie zonk

de mensenhandelsvloot.
De jager werd gedood

door een bosguerrilla
net als zijn collega-

vogelvrijenpesters.
Bereden houtvesters,

opgeleid tot ruiters,
vingen de uitbuiters.

Slimme bajesklanten
werden infiltranten

in ruil voor hun vrijheid.
Deze legereenheid

exploiteerde de biecht
om stil een sektesmiecht

in een ledig biechtbed
te doden als inzet.

Ze vingen de slavin,
die getuigde dan in

een open hoorzitting.
Er kwam dan vergeving

en voor de infiltrant
als voor de arrestant.

En er kwam dan meer hulp.
De kerk kroop uit zijn schulp.
Venters lengden hun drank
aan met aluin zo rank

dat de stad indommelt
als wordt opgetrommeld

de landheer zijn leger
als een stad-aanveger.

De landheers turfstekers
werden als grondbrekers

de afdeling genie.
Ze groeven tunnels, die

zijn soldaten brachten
binnen de stadsgrachten

en dan de stadsmuren,
door geen schot te vuren.

Die slachters en beukers
wisten wel wat leukers,
als droogstaande neukers,

te doen met de rotstad.
Alles ging dus daar plat.
Zodat niemand vergat,

wat men te wachten stond,
als een van hen ooit vond
dat het terecht bestond

en zou gaan herbouwen -
mensen en gebouwen -
met hernieuwd vertrouwen.

De bisschop sprak als wens
“Verloochen elke grens.
Laat geen steen, schrift of mens
in de stad overeind,
die de leer ondermijnt!”
De landheer ingeseind

ging toen helemaal los
en de richtte chaos.
Helaas bleek ook de klos

de zoon van de bisschop.
De zoon hing aan een strop
toen zijn vader reed op

zijn paard door de puinhoop.
Hij verloor elke hoop
op een goede afloop

Zeker, er zijn verschillende voorbeelden van steden en dorpen in Duitsland en Nederland die door menselijke activiteiten zijn verwoest. Een bekend voorbeeld in Duitsland is de stad Magdeburg, die in 1631 tijdens de Dertigjarige Oorlog vrijwel volledig werd verwoest door keizerlijke troepen. Dit werd bekend als het Bloedbad van Magdeburg en maakte een diepe indruk op de regio.
In Nederland is er het dorp Stavoren, dat ooit een welvarende handelsstad was aan de Zuiderzee. Volgens de sage van het Vrouwtje van Stavoren zou de stad zijn vervallen door de arrogantie en hebzucht van haar inwoners. Hoewel Stavoren niet volledig werd verwoest, verloor het wel zijn prominente positie als handelscentrum.
Deze verhalen kunnen dienen als inspiratie voor het lot van Pandemonium in je verhaal, waarbij menselijke daden leiden tot de ondergang van een eens florerende stad.

waailap = vlag met embleem bishop ter steun

De landheer had zijn leger opgetrommeld.
De schouw was ondertussen ingedommeld.

Hij had het er maar druk mee.
Al die vrouwen als trofee.

Toen hij zich uit eentje beroerde,
wist hij niet welk gevaar er loerde.

Totaal overmeesterd door de grote overmacht,
verloor de stad zijn soevereiniteit die nacht.

De landheer kon de haven wel waarderen.
Zijn soldaten mochten de slavinnen penetreren.

De priester en de burgemeester werden gedood.
Terwijl de landheer van de drie misdienaressen genoot.

De schouw werd voorgeleid
en hij was meteen bereid

om te vertellen was er met Emmy was gebeurd.
Als straf werd hij door jachthonden verscheurd.

De angstige stadjes verklaarden een voor een
hoe het was gegaan met Emmy en Heidi.

De soldaten dreven de mannen naar de kerk
en starten vervolgens een immens vuurwerk.

Ook de vrouwen werden door soldaten versleept.
Ze werden in de haven verscheept.

De misdienaressen vielen een ander lot ten deel.
Die werden toegevoegd aan de landheer’s bordeel.

Naast de kapel had ook de rest van de stad
inmiddels vuur en vlam gevat.

Doordat de schouw zo had geblunderd,
werd zijn gehele stad geplunderd.

De landheer was toch niet tevreden.
Zijn dochter had immers veel geleden.

De landheer wist dat Emmy zijn kind was.
Omdat hij dit in een brief van haar zus las.

Brieven die hij sinds haar vlucht onderschiep
terwijl hij met haar mammie sliep.

Emmy’s zuster werd uit het bordeel gehaald.
De landheer zette het haar vaar betaald.

Haar mère kon hij uiteindelijk vergeven
en is ook sindsdien uit het bordeel gebleven.

Zo leefde moeke, zuster en landheer samen
en moesten alle lijfeigenen hen gehoorzamen.

Gevormd door al het ondergane leed,
werden moeder en dochter net zo wreed.
